% Section 6 -- What minimal-data probes learn.
% 6.3 is the author's draft of 2026-09-05. Five figures in it conflicted with the committed results; all five were
% corrected to the results and the corrections were approved by the author on 2026-09-05, so the \disc flags are
% cleared. The refit passage was then rewritten by the author to state affirmatively that no evaluation data touched
% the probes -- a stronger fact than the original wording, and the one the corrected method actually supports.
\section{What Minimal-Data Probes Learn}
\label{sec:learn}

Section~\ref{sec:extraction} showed that where a probe is read changes what it appears to know. This section asks the
complementary question: with the extraction convention held fixed, what have these probes actually learned? The answer
differs by task, and in one case the more interesting finding is that the question cannot be answered with the
instruments available.

\subsection{Topic, not harm}
\label{sec:topicharm}

The harmful-content probes separate topic from not-topic rather than harmful from benign, and the two benign sets of
Section~\ref{sec:data} make this legible. Read at the final token of bare text, the shipped probes reach AUC
$0.90$--$1.00$ against XSTest---benign prompts that merely sound alarming---and only $0.46$--$0.67$ against
JailbreakBench benign prompts, which share their subject matter with harmful requests. The false-positive rates say
the same thing more bluntly: $0$--$17\,\%$ on XSTest against $30$--$94\,\%$ on the on-topic set. A detector that fires
on nine in ten benign requests about the same subject is not detecting harm.

This is what makes the pooled benign set a reporting hazard rather than a convenience. Pooling XSTest and
JailbreakBench yields an aggregate AUC of $0.81$--$0.91$, which reads as a working detector and conceals that
essentially all of its errors fall on one half of the benign data. It also explains the system-level result of
Section~\ref{sec:guard}: all $94$ of the probe pipeline's false positives are on-topic benign prompts.

Chat-templated reading substantially repairs this (Table~\ref{tab:extraction_convention}), which is the finding of
Section~\ref{sec:template}. What it does not do is make the probe a harm detector by construction, and the lexical
baselines of Table~\ref{tab:af_lexical} bound how much credit is due: the best cheap text feature reaches
$\aucstar{} = 0.710$, so the templated probe's margin over ``read the words'' is $+0.213$ to $+0.272$, not the whole
of its AUC.

\subsection{Paired contrastive training is not a training method}
\label{sec:paired}

A recurring proposal for minimal-data probes is to build the training set from matched pairs---each positive aligned
with a negative sharing topic, style and length---and to average the \emph{within-pair} differences, on the argument
that shared features cancel in the difference and only the safety-relevant direction survives.

The argument is sound and the method is not distinct. For $n$ complete, aligned pairs,
\begin{equation}
    \frac{1}{n}\sum_{i=1}^{n}\left(h^{+}_{i} - h^{-}_{i}\right)
    \;=\; \frac{1}{n}\sum_{i=1}^{n} h^{+}_{i} - \frac{1}{n}\sum_{i=1}^{n} h^{-}_{i},
    \label{eq:identity}
\end{equation}
which is exactly~\eqref{eq:direction}. The pairing changes the estimator not at all; the two agree to the last bit,
and our implementation asserts the equality on every model and mode rather than merely reporting it. Any difference
between a ``paired'' and an ``unpaired'' probe is therefore a difference in the training \emph{data}, and an
experiment comparing the two rules is measuring nothing.

That leaves the real question: does building the training set from on-topic pairs help? It is testable by holding the
estimator and the positive examples fixed and varying only the negatives, which is the \emph{off-topic} row of
Table~\ref{tab:extraction_convention}. Swapping $40$ on-topic benign twins for the off-topic negatives the shipped
probes were fitted with moves AUC against on-topic benign prompts by between $-0.003$ and $+0.018$---indistinguishable
from zero on every model---while the extraction change discussed above moves it by $+0.296$ to $+0.544$. The training
set was not where the leverage was.

We report this as a retirement rather than a refutation. The proposal is not wrong about pairing cancelling shared
features; it is wrong that this constitutes a distinct \emph{method}, and that part is settled by
equation~\eqref{eq:identity} rather than by measurement. What remains is the data construction, and there our claim is
the narrower empirical one: across seven models we observed no effect distinguishable from zero, against an extraction
change measured on the same probes that moved the same quantity by an order of magnitude more. We have not run an
equivalence test, so this bounds the effect as small on these models rather than establishing that it is nil.

\subsection{Policy compliance: when the evaluation cannot certify the claim}
\label{sec:apc}

Our third task asks whether a probe can distinguish giving advice from providing information---a speech-act
distinction underlying enterprise policies such as ``do not provide medical advice.'' The results here are of a
different kind from Section~\ref{sec:extraction}'s: we diagnose the failure, and then show that the natural instrument
for testing any repair is itself invalid at this text granularity.

\textbf{The minimal-data probe has no external validity.} A compliance probe trained in the standard minimal-data
recipe---fifteen authored contrastive pairs per policy, evaluated on six held-out pairs---achieves
a perfect held-out score. On an external evaluation set of $360$ labelled responses it collapses: AUC $0.14$--$0.58$
across models (Table~\ref{tab:apc_external}), and at its own calibrated threshold it flags every case---every advice
response, every information response, and all twenty refusals---in four of the six model/mode cells
(Table~\ref{tab:apc_refusal}). The perfect original score is an artifact of a six-example split; the refusal result
illustrates a measurement error we suggest is general, in which a probe's detections are never separated from the
model's own refusal behaviour, conflating ``the probe caught it'' with ``the model was never going to comply.''

% GENERATED by paper/tools/gen_tables.py -- do not edit by hand.
% sources: results/gates/b3/*.json
\begin{table*}[!t]
\caption{Policy-compliance probes on the external evaluation set, chat-templated read. \emph{shipped} is the
minimal-data probe as distributed; \emph{fresh} is the same recipe refitted on the same authored training pairs
at the same layer --- neither is fitted to the external set. Refusals are excluded from the AUC and reported
separately in Table~\ref{tab:apc_refusal}.}
\label{tab:apc_external}
\centering
\small
\begin{tabular}{llcccccc}
\toprule
Model & Probe & medical & medical$^{*}$ & financial & legal & flag adv. & flag inf. \\
\midrule
Gemma-2-9B & shipped & 0.894 & 0.893 & 0.638 & 0.617 & 1.00 & 1.00 \\
 & fresh & 0.867 & 0.889 & 0.962 & 1.000 & 0.75 & 0.15 \\
Llama-3.1-8B & shipped & 0.144 & 0.126 & 0.270 & 0.128 & 0.38 & 0.94 \\
 & fresh & 0.991 & 0.996 & 1.000 & 0.997 & 0.89 & 0.02 \\
Llama-3.2-3B & shipped & 0.408 & 0.395 & 0.598 & 0.503 & 1.00 & 1.00 \\
 & fresh & 0.986 & 0.994 & 0.978 & 1.000 & 0.95 & 0.05 \\
\bottomrule
\end{tabular}
\end{table*}
% GENERATED by paper/tools/gen_tables.py -- do not edit by hand.
% sources: results/gates/b3/*.json
\begin{table*}[!t]
\caption{The refusal band, reported separately because it is invisible inside an advice-versus-information
AUC that excludes it. A compliance probe that flags refusals fires on exactly the responses the policy is
trying to elicit. The last column is the probe's AUC separating refusals from compliant information; $0.5$
means it cannot tell them apart.}
\label{tab:apc_refusal}
\centering
\small
\begin{tabular}{lllcccc}
\toprule
Model & Mode & Probe & $n$ & flag rate & mean score & AUC vs inf. \\
\midrule
Gemma-2-9B & raw & shipped & 20 & 1.00 & +0.016 & 0.312 \\
Gemma-2-9B & raw & fresh & 20 & 0.00 & -0.208 & 0.879 \\
Gemma-2-9B & templated & shipped & 20 & 1.00 & +0.002 & 0.203 \\
Gemma-2-9B & templated & fresh & 20 & 0.60 & -0.038 & 0.167 \\
Llama-3.1-8B & raw & shipped & 20 & 0.40 & -0.005 & 0.680 \\
Llama-3.1-8B & raw & fresh & 20 & 0.00 & -0.064 & 0.716 \\
Llama-3.1-8B & templated & shipped & 20 & 0.15 & -0.014 & 0.967 \\
Llama-3.1-8B & templated & fresh & 20 & 0.70 & -0.056 & 0.057 \\
Llama-3.2-3B & raw & shipped & 20 & 1.00 & +0.006 & 0.327 \\
Llama-3.2-3B & raw & fresh & 20 & 0.00 & -0.090 & 0.916 \\
Llama-3.2-3B & templated & shipped & 20 & 1.00 & -0.004 & 0.918 \\
Llama-3.2-3B & templated & fresh & 20 & 0.50 & -0.019 & 0.102 \\
\bottomrule
\end{tabular}

\footnotesize The shipped probe flags every case in its evaluation cell --- all advice, all information and
all 20 refusals --- in 4 of the $6$ model/mode cells.
\end{table*}

\textbf{Refit probes pass---and the pass is uninterpretable.} Probes refit by the same recipe on the same fifteen
authored pairs---differing from the shipped probes only in extraction convention---reach AUC $0.99$ within-topic and
$0.98$--$1.00$ cross-topic on the external set. No evaluation data touched the probes; apparently strong evidence of
speech-act reading. But a two-token regex---the presence of second-person pronouns---achieves $0.93$--$0.95$ on the same contrasts, and $1.000$
on the hardest cells (Table~\ref{tab:apc_lexical}). The evaluation cannot distinguish a semantic classifier from a
lexical-cue detector, because in this authored set the cue and the construct coincide.

% GENERATED by paper/tools/gen_tables.py -- do not edit by hand.
% sources: results/gates/lexical_baselines.json
\begin{table*}[!t]
\caption{Lexical baselines on the policy-compliance contrasts: no model, no activations, just a count. On the
v1 evaluation set a second-person pronoun count matches or beats the probe, so that set cannot distinguish a
speech-act classifier from a cue detector. The v2 rows are the instrument built to break the cue.}
\label{tab:apc_lexical}
\centering
\small
\begin{tabular}{llccc c}
\toprule
Set & Scope & $n$ & second person & modals & length \\
\midrule
v1 medical & all bands & 220 & \textbf{0.936} & 0.885 & 0.514 \\
v1 medical & hard bands & 110 & \textbf{0.955} & 0.862 & 0.524 \\
v1 financial & all bands & 56 & \textbf{0.929} & 0.801 & 0.500 \\
v1 financial & hard bands & 28 & \textbf{0.929} & 0.768 & 0.513 \\
v1 legal & all bands & 56 & \textbf{0.946} & 0.824 & 0.518 \\
v1 legal & hard bands & 28 & \textbf{1.000} & 0.906 & 0.546 \\
v2 & all cells & 80 & \textbf{0.513} & 0.500 & 0.512 \\
v2 & incongruent cells only & 40 & \textbf{0.000} & 0.500 & 0.501 \\
\bottomrule
\end{tabular}
\end{table*}

\textbf{Controlling the cue dissolves the construct.} We attempted to build the discriminating instrument: an
evaluation set in which lexical cues and the advice/information label are decorrelated by construction. Three
successive drafts failed in three distinct ways (Table~\ref{tab:apc_construction}). The first controlled second-person
marking; a directive-modal regex then separated the classes at $0.875$. The second controlled modals; a bag-of-words
classifier separated them at $1.000$ via register---professional nouns such as \emph{clinicians} and \emph{guidelines} appearing in $28$ of $40$ information cases against $8$ of $40$
advice cases. The third drove every measured lexical cue to chance---pronouns $0.544$, modals $0.545$,
professional-register nouns $0.500$, length $0.485$, omnibus lexical classifier $0.478$ against a within-pair-swap
null of $0.435 \pm 0.084$---and in doing so dissolved the labels themselves: median within-pair similarity after stopword removal reached $0.814$,
and members of the ``advice'' class no longer carried directive force under our own pre-registered labelling rule.

% GENERATED by paper/tools/gen_tables.py -- do not edit by hand.
% sources: results/gates/validity_checks.json
\begin{table*}[!t]
\caption{Three attempts to build an evaluation set in which lexical cues and the advice/information label are
decorrelated. Each draft controlled the cue the previous one exposed, and each exposed a new one; the third
drove every measured cue to chance and dissolved the labels instead.}
\label{tab:apc_construction}
\centering
\small
\begin{tabular}{c p{0.15\textwidth} p{0.40\textwidth} p{0.26\textwidth}}
\toprule
Draft & Cue controlled & Cue exposed & Value \\
\midrule
1 & second-person pronouns & directive modals & 0.875 (regex AUC) \\
2 & directive modals & professional/third-party register (clinicians, guidelines, pharmacists: 28/40 information cases vs 8/40 advice) & 1.000 (omnibus LOO) \\
3 & register, via 40 minimal pairs & label validity dissolved & 0.814 (median within-pair content similarity after removing pronouns, articles and copulas) \\
\bottomrule
\end{tabular}

\footnotesize Draft~3 measured cues: second person $0.544$, directive modals $0.545$, professional-register nouns $0.500$, character length $0.485$; omnibus bag-of-words under leave-pair-out cross-validation $0.478$ against a within-pair-swap null of $0.435 \pm 0.084$. Label validity: median within-pair content similarity $0.814$, with 10 of $40$ pairs above $0.90$.
\end{table*}

The omnibus figure needs one methodological note, developed in Section~\ref{sec:leakage}: on minimal-pair data
leave-one-out cross-validation is invalid, because a held-out item's twin sits in the training fold carrying the
opposite label and the classifier reliably anti-predicts. Measured that way draft~3 scores $0.019$---near-perfect
separation with the sign inverted---which says nothing about lexical separability. Leave-pair-out against a
within-pair-swap null is the correct procedure and is what the figures above report.

\textbf{The finding.} The three failures share one cause, which we state as a result rather than as a limitation: in
short text, directive force is substantially constituted by its lexical markers. Pronoun address, deontic modality,
and register are not confounds of the advice/information distinction at this granularity; they are much of what the
distinction consists in. It follows that short-text authored evaluations cannot certify a semantic claim for any
compliance classifier---probe, fine-tuned model, or otherwise---because any instrument that removes the cues removes
the construct.

Two practical consequences. First, for short responses, a lexical classifier may simply be an adequate compliance
control, and claims that a representation-level method outperforms one require evidence such evaluations cannot
provide. Second, certifying semantic speech-act classification requires naturally occurring multi-sentence responses,
where directive force is carried structurally across discourse---an evaluation we scope as future work rather than
attempt with an instrument we have shown cannot answer the question.

We are careful about what this does and does not say. It is not a claim that these probes cannot read speech acts. It
is the narrower and better-supported claim that no evidence we can construct at this text granularity could establish
that they do.
